%----------------------------------------------------------%
% !TEX outputDirectory = build

\documentclass[reqno,a4paper,12pt]{amsart}
\def\thmcolor{Sepia!5!white}
\def\lemcolor{NavyBlue!5!white}
\def\propcolor{PineGreen!5!white}
\def\corcolor{Orange!5!white}
\def\clmcolor{Cyan!5!white}
\def\quecolor{ForestGreen!5!white}
\def\defcolor{OliveGreen!5!white}
\def\obscolor{Plum!5!white}
\def\remcolor{black!5!white}
\input{header.tex}

\usepackage{tikz}
\usetikzlibrary{patterns}
\usepackage{algorithm}
\usepackage[noend]{algpseudocode}

% Temporary nonsense
\usepackage{pgfplots}
\pgfplotsset{compat=1.15,width=10cm,height=10cm}

%----------------------------------------------------------%
%% Document specific definitions

%% Comments
\newcommand{\victor}[1]{\textcolor{magenta}{\textbf{[#1]}}}
\newcommand{\afonso}[1]{\textcolor{red}{\textbf{[#1]}}}
\newcommand{\roberto}[1]{\textcolor{blue}{\textbf{[#1]}}}
\newcommand{\felipe}[1]{\textcolor{orange}{\textbf{[#1]}}}

%----------------------------------------------------------%

\begin{document}

\title{The number of strong Sidon sets}

\author{Felipe Guimarães}
\author{Roberto Parente}
\author{Afonso Sant'Anna}
\author{Victor Souza}

%----------------------------------------------------------%


\maketitle

%----------------------------------------------------------%

\section{Sketch introduction}

In enumerative extremal combinatorics, there are several results that say that the number of structures without a substructures satisfy the asymptotic
\begin{equation}
    \label{eq:extremal-counting}
    \card{\set{\text{structures}}} = 2^{\text{extremal number} + o(1)}.
\end{equation}
Example include sum-free sets, $H$-free graphs, whatever [CITE].
There are even examples when \eqref{eq:extremal-counting} is known to hold when the order of growth of the extremal number $\text{extremal number}$ is not know, for instance for $k-AP$s [CITE].

A lower bound of $2^{\text{extremal number}}$ is always available for $\card{\set{\text{structures}}}$ from monotonicity.
The difficulty lies in proving the corresponding upper bound, which more or less correspond to showing that `most' structures are obtained in this way.
This is not quite true as the $2^{o(1)}$ term can be really massive, \eqref{eq:extremal-counting} says that this is true in a logarithmic sense.

However, even for monotone properties, \eqref{eq:extremal-counting} is not always true.
A famous example where this is not true is in the count of Sidon sets, namely subsets $A \subseteq [N]$ with definition blabla.
The best known results are the following
\begin{equation}
    \label{eq:sidon-bounds}
    (2.1)^{\sqrt{n}} \leq \card{\set{\text{Sidon sets}}} \leq 10^{\sqrt{n}}.
\end{equation}

Two features are striking from \eqref{eq:sidon-bounds}: that we do not known the correct constant $C$ such that we have $C^{\sqrt{n} + o(1)}$ Sidon sets and that this constant is not $2$.
Indeed, we have $C > 2$ from a construction of Saxton and Thomason,
In fact, not a lot is know about the structure of large Sidon sets, cite some papers.
This is a serious obstacle to obtain a corresponding lower bound.
The upper bounds are due to BOB, improving on an earlier bound of $1000^{\sqrt{n}}$ from BEB.

We consider the analogous enumeration problem for $k$-\emph{strong} Sidon sets, namely, subsets $A \subseteq [N]$ with the property that for every $a,b,c,d \in A$ with $\set{a,b} \cap \set{c,d} = \emptyset$, we have $\abs{(a + b) - (c + d)} \geq k$.
These were considered by Yoshi et al in [CITE].
They showed that

%--------------------------------------%
\begin{theorem}
    \label{thm:disjointly disjointly $k$-Sidon-extremal}
    If $A \subseteq [n]$ is a $k$-strong Sidon set, then $\card{A} \leq \sqrt{n/k}$.
    This bound can be attained.
\end{theorem}
%--------------------------------------%

As they only consider the case $k = n^{\alpha}$ in their work, we include a full proof.
However, our main contribution is to prove the following analogue of \eqref{eq:sidon-bounds}.

%--------------------------------------%
\begin{theorem}
    \label{thm:disjointly disjointly $k$-Sidon-enumeration}
    There are explicit constants $1 < c \leq C$ such that
    \begin{equation*}
        2^{c \log k \sqrt{n/k}} \leq \card{\set{\text{$k$-strong Sidon sets}}} \leq 2^{C \log k \sqrt{n/k}},
    \end{equation*}
    uniformly for all $1 \leq k \leq n/(\log n)^8$.
\end{theorem}
%--------------------------------------%

In view of the extremal number from \Cref{thm:disjointly disjointly $k$-Sidon-extremal}, the bounds in \Cref{thm:disjointly disjointly $k$-Sidon-enumeration} furnishes an infinite family of combinatorial structures for which the heuristics in \eqref{eq:extremal-counting} is not valid.
More interestingly, we have been able to determine the dependency in $k$.

The extra factor of $\log k$ in the enumeration can be interpreted as saying that when choosing a typical $k$-strong Sidon set, each element contributes on average with $\Theta(k)$ choices.
Even when $k = 1$, this contributes with a nontrivial gain in entropy.

As many results of it's type, the upper bound in \Cref{thm:disjointly disjointly $k$-Sidon-enumeration} uses the graph container theorem [CITE] and a novel supersaturation result (Reference).
The lower bound is a direct generalisation of the Saxton and Thomason construction.



%----------------------------------------------------------%

\section{Sketch}

$[N] \defined \set{1, \dotsc, N}$.

\begin{definition}[$k$-strong Sidon set]
    \label{def:k-strong-sidon}
    A subset $A\subseteq[N]$ is called \emph{$k$-strong Sidon} if, for
    every $x,y,z,w\in A$ satisfying
    \[
        x<y\leq z<w,
    \]
    we have
    \[
        \abs{(x+w)-(y+z)}\geq k.
    \]
\end{definition}

%--------------------------------------%
\begin{definition}[Disjointly $k$-Sidon set]
    \label{def:disjoint-disjointly disjointly $k$-Sidon}
    A subset $A\subseteq[N]$ is called \emph{disjointly $k$-Sidon} if,
    for every $a,b,c,d\in A$ satisfying
    \[
        \set{a,b}\cap\set{c,d}=\emptyset,
    \]
    we have
    \[
        \abs{(a+b)-(c+d)}\geq k.
    \]

    Repetitions within each pair are allowed; that is, we may have
    $a=b$ or $c=d$. The condition only requires that no element occurring
    in the first pair occurs in the second pair.
\end{definition}

\begin{definition}
    A subset $A \subseteq [N]$ is $\delta$-separated if for every $a,b \in A$ with $a \neq b$, we have $\abs{a - b} \geq \delta$.
\end{definition}

The definition of a \(k\)-strong Sidon set imposes a condition only on
quadruples appearing in the order
\[
    x<y\leq z<w,
\]
and at first sight this says nothing about how close together two
elements of \(A\) are allowed to be. As we now show, this single
condition already forces strong restrictions on the small gaps of
\(A\).

Suppose that \(A\) contains two disjoint pairs
\[
    \{i,\,i+\delta\}
    \qquad\text{and}\qquad
    \{j,\,j+\alpha\},
    \qquad
    0<\delta<k,
    \quad
    0<\alpha<k,
\]
both violating \(k\)-separation. Exchanging the two pairs if
necessary, we may assume \(i<j\). It remains only to locate
\(i+\delta\) relative to \(j\) and \(j+\alpha\), which leaves exactly
three possibilities:
\[
    i+\delta<j,
    \qquad
    j<i+\delta<j+\alpha,
    \qquad
    j+\alpha<i+\delta,
\]
corresponding to the two pairs being \emph{separated}, \emph{crossing},
or \emph{nested}, respectively. We rule out each case using nothing
beyond Definition~\ref{def:k-strong-sidon}.

First, suppose that the two pairs are separated, so that the four
elements occur in the order
\[
    i<i+\delta<j<j+\alpha.
\]
Applying the \(k\)-strong Sidon condition to this ordered quadruple
gives
\[
    \abs{\bigl(i+(j+\alpha)\bigr)-\bigl((i+\delta)+j\bigr)}
    =
    \abs{\alpha-\delta}
    \geq k.
\]
However, since \(0<\alpha,\delta<k\), we have
\(-k<\alpha-\delta<k\), and hence \(\abs{\alpha-\delta}<k\), a
contradiction.

Suppose next that the two pairs are nested, so that
\[
    i<j<j+\alpha<i+\delta.
\]
Write \(j=i+\gamma\) and \(j+\alpha=i+\beta\), so that
\(0<\gamma<\beta<\delta<k\). The four elements then occur in the
order
\[
    i<i+\gamma<i+\beta<i+\delta,
\]
and the \(k\)-strong Sidon condition gives
\[
    \abs{\bigl(i+(i+\delta)\bigr)-\bigl((i+\gamma)+(i+\beta)\bigr)}
    =
    \abs{\delta-(\gamma+\beta)}
    \geq k.
\]
On the other hand, \(0<\gamma+\beta<2\delta\), so
\[
    -\delta<\delta-(\gamma+\beta)<\delta,
\]
and thus
\[
    \abs{\delta-(\gamma+\beta)}<\delta<k,
\]
again a contradiction.

Finally, suppose that the two pairs cross, so that
\[
    i<j<i+\delta<j+\alpha.
\]
Applying the \(k\)-strong Sidon condition to this ordered quadruple
yields
\[
    \abs{\bigl(i+(j+\alpha)\bigr)-\bigl(j+(i+\delta)\bigr)}
    =
    \abs{\alpha-\delta}
    \geq k,
\]
which contradicts \(\abs{\alpha-\delta}<k\) exactly as in the
separated case.

Thus, two disjoint pairs of elements of a \(k\)-strong Sidon set
cannot both have distance less than \(k\). The following lemma
records the underlying four-point phenomenon in a form that will be
useful throughout the proof.

%--------------------------------------%
\begin{lemma}[Four-point pairing lemma]
    \label{lem:four-point-pairing}
    Let \(A\subseteq[N]\) be a \(k\)-strong Sidon set, and let
    \[
        x_1<x_2<x_3<x_4
    \]
    be four elements of \(A\). Then every partition of
    \[
        \set{x_1,x_2,x_3,x_4}
    \]
    into two pairs gives pair sums whose difference has absolute value
    at least \(k\).
\end{lemma}
%--------------------------------------%

\begin{proof}
    Write
    \[
        p\defined x_2-x_1,
        \qquad
        q\defined x_3-x_2,
        \qquad
        r\defined x_4-x_3.
    \]
    Since
    \[
        x_1<x_2\leq x_3<x_4,
    \]
    the \(k\)-strong Sidon property gives
    \[
        \abs{
            (x_1+x_4)-(x_2+x_3)
        }
        \geq k.
    \]
    Moreover,
    \[
        (x_1+x_4)-(x_2+x_3)
        =
        r-p,
    \]
    and hence
    \[
        \abs{r-p}\geq k.
    \]

    There are exactly three partitions of four elements into two
    unordered pairs. Besides
    \[
        \set{x_1,x_4}
        \mid
        \set{x_2,x_3},
    \]
    the other two are
    \[
        \set{x_1,x_3}
        \mid
        \set{x_2,x_4}
    \]
    and
    \[
        \set{x_1,x_2}
        \mid
        \set{x_3,x_4}.
    \]

    For the second partition,
    \[
        \abs{
            (x_1+x_3)-(x_2+x_4)
        }
        =
        p+r.
    \]
    We claim that
    \[
        \abs{r-p}<p+r.
    \]
    Indeed, if \(r\geq p\), then
    \[
        \abs{r-p}=r-p,
    \]
    and
    \[
        (p+r)-\abs{r-p}
        =
        (p+r)-(r-p)
        =
        2p
        >
        0.
    \]
    If instead \(r<p\), then
    \[
        \abs{r-p}=p-r,
    \]
    and
    \[
        (p+r)-\abs{r-p}
        =
        (p+r)-(p-r)
        =
        2r
        >
        0.
    \]
    Thus, in either case,
    \[
        \abs{r-p}<p+r.
    \]
    Since
    \[
        \abs{r-p}\geq k,
    \]
    it follows that
    \[
        p+r\geq k.
    \]

    For the third partition,
    \[
        \abs{
            (x_1+x_2)-(x_3+x_4)
        }
        =
        p+2q+r.
    \]
    Since \(q>0\),
    \[
        p+r<p+2q+r.
    \]
    Consequently,
    \[
        k
        \leq
        \abs{r-p}
        <
        p+r
        <
        p+2q+r,
    \]
    and therefore this third difference is also at least \(k\).

    Hence every partition of
    \(\set{x_1,x_2,x_3,x_4}\) into two pairs gives pair sums differing
    by at least \(k\).
\end{proof}
%--------------------------------------%

The four-point pairing lemma implies that all failures of
\(k\)-separation inside a \(k\)-strong Sidon set must be concentrated
around a single element. This is the content of the next lemma.

%--------------------------------------%
\begin{lemma}[\(k\)-separated core]
    \label{lem:k-separated-core}
    Let \(A\subseteq[N]\) be a \(k\)-strong Sidon set. Then there exists
    \(A'\subseteq A\) such that
    \[
        \card{A'}\geq\card{A}-1
    \]
    and \(A'\) is \(k\)-separated.
\end{lemma}
%--------------------------------------%

\begin{proof}
    Define a graph \(G\) with vertex set \(A\) by declaring two distinct
    elements \(a,b\in A\) adjacent whenever
    \[
        \abs{a-b}<k.
    \]
    Thus, the edges of \(G\) record precisely the pairs that violate
    \(k\)-separation.

    We first show that \(G\) contains no two disjoint edges. Suppose,
    for a contradiction, that
    \[
        \set{a,b}
        \qquad\text{and}\qquad
        \set{c,d}
    \]
    are disjoint edges. Relabelling inside each pair if necessary,
    assume
    \[
        a<b
        \qquad\text{and}\qquad
        c<d.
    \]
    Since both pairs are edges,
    \[
        0<b-a<k
        \qquad\text{and}\qquad
        0<d-c<k.
    \]
    Therefore,
    \[
        \abs{(b-a)-(d-c)}<k.
    \]
    But
    \[
        (b-a)-(d-c)
        =
        (b+c)-(a+d).
    \]
    Since \(a,b,c,d\) are four distinct elements, the pairs
    \[
        \set{b,c}
        \qquad\text{and}\qquad
        \set{a,d}
    \]
    form a partition of these four elements. By
    Lemma~\ref{lem:four-point-pairing},
    \[
        \abs{(b+c)-(a+d)}
        \geq k,
    \]
    a contradiction.

    We next show that \(G\) contains no triangle. Suppose that
    \[
        x<y<z
    \]
    span a triangle. In particular,
    \[
        xz\in E(G),
    \]
    so
    \[
        z-x<k.
    \]
    Since
    \[
        x<y\leq y<z,
    \]
    the \(k\)-strong Sidon property gives
    \[
        \abs{x+z-2y}\geq k.
    \]
    On the other hand,
    \[
        \abs{x+z-2y}
        =
        \abs{(y-x)-(z-y)}.
    \]
    Since both \(y-x\) and \(z-y\) are positive,
    \[
        \abs{(y-x)-(z-y)}
        \leq
        (y-x)+(z-y)
        =
        z-x
        <
        k,
    \]
    a contradiction.

    We now determine the structure of \(G\).
    If \(G\) has no edges, then \(A\) itself is \(k\)-separated and we
    may take
    \[
        A'\defined A.
    \]

    Suppose that \(G\) has at least one edge. If it has exactly one
    edge, choose either endpoint of that edge and call it \(v\).

    Otherwise, choose two distinct edges of \(G\). Since \(G\) has no
    two disjoint edges, they must intersect. Write them as
    \[
        uv
        \qquad\text{and}\qquad
        uw.
    \]
    We claim that every edge of \(G\) contains \(u\).

    Indeed, let \(e\) be any edge. Since \(G\) has no two disjoint
    edges, \(e\) must intersect both \(uv\) and \(uw\). If \(u\notin e\),
    then the only way for \(e\) to meet both of these edges is
    \[
        e=vw.
    \]
    But then
    \[
        uv,\quad uw,\quad vw
    \]
    form a triangle, contradicting what we proved above.

    Hence all edges of \(G\) share a common vertex \(u\). Set
    \[
        A'\defined A\setminus\set{u}.
    \]
    Then \(G[A']\) has no edges. Therefore, for every pair of distinct
    \(a,b\in A'\),
    \[
        \abs{a-b}\geq k.
    \]
    Thus \(A'\) is \(k\)-separated, and clearly
    \[
        \card{A'}
        \geq
        \card{A}-1.
    \]
\end{proof}
%--------------------------------------%


%--------------------------------------%
\begin{definition}
    \label{def:auxiliary-multigraph}
    Given \(A,U\subseteq[N]\), define the loopless multigraph
    \[
        \cH_k(A,U)
    \]
    as follows. Its vertex set is \(A\). For every pair of distinct
    vertices \(a,a'\in A\), with \(a<a'\), and every ordered pair
    \[
        (u,u')\in U^2,
        \qquad
        u\neq u',
    \]
    we place one edge between \(a\) and \(a'\) whenever
    \[
        \abs{(a+u)-(a'+u')}<k.
    \]
    We say that this edge has \emph{type} \((u,u')\). Thus, parallel
    edges correspond to distinct ordered pairs \((u,u')\) satisfying
    the inequality above.
\end{definition}
%--------------------------------------%

%--------------------------------------%
\begin{proposition}
    \label{prop:sidon-independent}
    Let
    \[
        I\subseteq A\subseteq[N]\setminus U.
    \]
    If \(U\cup I\) is \(k\)-strong Sidon, then \(I\) is an independent
    set of \(\cH_k(A,U)\).
\end{proposition}
%--------------------------------------%

\begin{proof}
    Suppose, to the contrary, that \(I\) is not independent. Then there
    exist distinct vertices \(a,a'\in I\) and distinct elements
    \(u,u'\in U\) such that
    \[
        \abs{(a+u)-(a'+u')}<k.
    \]
    Since
    \[
        A\cap U=\emptyset,
    \]
    the four elements
    \[
        a,a',u,u'
    \]
    are pairwise distinct.

    Since \(U\cup I\) is \(k\)-strong Sidon,
    Lemma~\ref{lem:four-point-pairing}, applied to the partition
    \[
        \set{a,u}
        \mid
        \set{a',u'},
    \]
    gives
    \[
        \abs{(a+u)-(a'+u')}\geq k,
    \]
    a contradiction.
\end{proof}
%--------------------------------------%

We next control the multiplicity of the auxiliary multigraph.

%--------------------------------------%
\begin{proposition}
    \label{prop:multiplicity}
    Let
    \[
        k\geq\Delta\geq1.
    \]
    Suppose that \(U\subseteq[N]\) is a \(\Delta\)-separated
    \(k\)-strong Sidon set and that
    \[
        A\subseteq[N]\setminus U.
    \]
    Then every edge of \(\cH_k(A,U)\) has multiplicity at most
    \[
        \left\lceil\frac{2k-1}{\Delta}\right\rceil.
    \]
\end{proposition}
%--------------------------------------%

\begin{proof}
    Fix distinct vertices \(a,a'\in A\), with \(a<a'\), and suppose
    that the edge \(aa'\) has multiplicity \(m\). Then there exist
    \(m\) distinct ordered pairs
    \[
        (u_1,u'_1),\dotsc,(u_m,u'_m)\in U^2,
        \qquad
        u_i\neq u'_i,
    \]
    such that
    \[
        \abs{(a+u_i)-(a'+u'_i)}<k
    \]
    for every \(i\in[m]\).

    Define
    \[
        d_i\defined u_i-u'_i.
    \]
    Since
    \[
        (a+u_i)-(a'+u'_i)
        =
        d_i-(a'-a),
    \]
    we have
    \[
        \abs{d_i-(a'-a)}<k.
    \]
    As all quantities are integers,
    \[
        a'-a-(k-1)
        \leq
        d_i
        \leq
        a'-a+(k-1).
    \]
    Thus all the \(d_i\)'s lie among \(2k-1\) consecutive integers,
    and hence in an interval of diameter \(2k-2\).

    We claim that the \(d_i\)'s are pairwise
    \(\Delta\)-separated. Suppose, to the contrary, that for some
    distinct \(i,j\in[m]\),
    \[
        \abs{d_i-d_j}<\Delta.
    \]
    Since
    \[
        d_i-d_j
        =
        (u_i-u'_i)-(u_j-u'_j)
        =
        (u_i+u'_j)-(u_j+u'_i),
    \]
    it follows that
    \begin{equation}
        \label{eq:multiplicity-close-sums}
        \abs{(u_i+u'_j)-(u_j+u'_i)}
        <
        \Delta
        \leq k.
    \end{equation}

    We distinguish the possible coincidences among
    \[
        u_i,\quad u'_i,\quad u_j,\quad u'_j.
    \]

    If these four elements are pairwise distinct, then the pairs
    \[
        \set{u_i,u'_j}
        \qquad\text{and}\qquad
        \set{u_j,u'_i}
    \]
    form a partition of four distinct elements of \(U\). By
    Lemma~\ref{lem:four-point-pairing},
    \[
        \abs{(u_i+u'_j)-(u_j+u'_i)}
        \geq k,
    \]
    contradicting \eqref{eq:multiplicity-close-sums}.

    Suppose next that
    \[
        u_i=u_j.
    \]
    Since the ordered pairs
    \((u_i,u'_i)\) and \((u_j,u'_j)\) are distinct, we must have
    \[
        u'_i\neq u'_j.
    \]
    Therefore,
    \[
        \abs{d_i-d_j}
        =
        \abs{u'_i-u'_j}
        \geq\Delta,
    \]
    since \(U\) is \(\Delta\)-separated, again a contradiction.

    Similarly, if
    \[
        u'_i=u'_j,
    \]
    then \(u_i\neq u_j\), and
    \[
        \abs{d_i-d_j}
        =
        \abs{u_i-u_j}
        \geq\Delta.
    \]

    It remains to consider the two possible cross-identifications.
    Suppose first that
    \[
        u_i=u'_j.
    \]
    Write
    \[
        x\defined u_i=u'_j,
        \qquad
        y\defined u'_i,
        \qquad
        z\defined u_j.
    \]
    Then
    \[
        d_i-d_j
        =
        2x-y-z.
    \]

    If \(y=z\), then
    \[
        \abs{d_i-d_j}
        =
        2\abs{x-y}
        \geq
        2\Delta
        \geq
        \Delta.
    \]

    Suppose now that \(x,y,z\) are pairwise distinct. If \(x\) lies
    between \(y\) and \(z\), then, after possibly interchanging \(y\)
    and \(z\), we may write
    \[
        y<x<z.
    \]
    Since
    \[
        y<x\leq x<z,
    \]
    the \(k\)-strong Sidon property gives
    \[
        \abs{y+z-2x}
        \geq k
        \geq\Delta.
    \]
    Hence
    \[
        \abs{d_i-d_j}
        =
        \abs{2x-y-z}
        \geq\Delta.
    \]

    If \(x\) does not lie between \(y\) and \(z\), then \(x\) lies
    strictly to one side of both. For instance, if
    \[
        x<y,z,
    \]
    then
    \[
        \abs{2x-y-z}
        =
        \abs{x-y}+\abs{x-z}
        \geq
        \Delta.
    \]
    The case \(x>y,z\) is analogous. Thus in every case
    \[
        \abs{d_i-d_j}\geq\Delta,
    \]
    a contradiction.

    Finally, suppose that
    \[
        u'_i=u_j.
    \]
    This case is symmetric to the previous one and again gives
    \[
        \abs{d_i-d_j}\geq\Delta.
    \]

    Hence the \(d_i\)'s are pairwise \(\Delta\)-separated.

    Since they lie in an interval of diameter \(2k-2\), we must have
    \[
        (m-1)\Delta\leq2k-2.
    \]
    Consequently,
    \[
        m
        \leq
        1+\floor{\frac{2k-2}{\Delta}}
        =
        \left\lceil\frac{2k-1}{\Delta}\right\rceil.
    \]
\end{proof}
%--------------------------------------%

\afonso{We can simplify the proof of Proposition~\ref{prop:multiplicity}
    by first proving the following auxiliary lemma: if \(U\) is a
    \(\Delta\)-separated \(k\)-strong Sidon set, with
    \(\Delta\leq k\), then the pair-sum multiset \(U\oplus U\) is
    \(\Delta\)-separated. Indeed, for two distinct witnesses
    \((u_i,u_i')\) and \((u_j,u_j')\),
    \[
        d_i-d_j
        =
        (u_i+u_j')-(u_j+u_i'),
    \]
    and the unordered pairs
    \[
        \set{u_i,u_j'}
        \qquad\text{and}\qquad
        \set{u_j,u_i'}
    \]
    are distinct as multisets. The \(\Delta\)-separation of
    \(U\oplus U\) would then imply immediately that
    \[
        \abs{d_i-d_j}\geq\Delta,
    \]
    avoiding the case analysis concerning possible coincidences among
    \(u_i,u_i',u_j,u_j'\).}

The \(k\)-separated core provided by
Lemma~\ref{lem:k-separated-core} gives the following form of the
multiplicity estimate that will be used below.

%--------------------------------------%
\begin{corollary}
    \label{cor:multiplicity}
    Suppose that \(U\subseteq[N]\) is a \(k\)-separated
    \(k\)-strong Sidon set and
    \[
        A\subseteq[N]\setminus U.
    \]
    Then every edge of \(\cH_k(A,U)\) has multiplicity at most \(2\).
\end{corollary}
%--------------------------------------%

\begin{proof}
    Apply Proposition~\ref{prop:multiplicity} with
    \[
        \Delta=k.
    \]
    Then every edge has multiplicity at most
    \[
        \left\lceil\frac{2k-1}{k}\right\rceil
        \leq2.
    \]
\end{proof}

\begin{lemma}[Supersaturation]
    \label{lem:supersaturation}
    For any $k \geq 1$ and $N \geq k$, we have that
    \begin{equation*}
        e(\cH_k(A,U)) \geq \frac{k \card{A}^2 \card{U}^2}{8N}
    \end{equation*}
    for every subsets $A, U \subseteq [N]$ satisfying $\card{A}\card{U} \geq 60 N$.
\end{lemma}
%--------------------------------------%
\begin{proof}
    Let $M \defined \card{A}\card{U}$ and $N' \defined 2N + k - 1$.

    Consider the auxiliary bipartite graph $F$ with disjoint parts $A$ and $[2N]$ where $a \in A$ is adjacent to $m \in [2N]$ if there exists $u \in U$ with $m = a + u$.
    Thus, each vertex $a \in A$ has precisely $\card{U}$ neighbours in $[2N]$, namely, the set of neighbours of $a$ is $a + U$.
    Therefore $e(F) = \card{A} \card{U} = M$.

    Let $\cX = \set{x \in \ZZ \st -k+1 \leq x \leq 2N - 1}$ and note that $\card{\cX} = N'$.
    For each $x \in \cX$, consider the interval
    \begin{equation*}
        I_x \defined \set{x + 1, \dotsc, x + k} = x + [k].
    \end{equation*}
    Given subsets $S \subseteq A$ and $T \subseteq [2N]$, we denote by $e_F(S,T)$ the number of edges in $F$ which are incident to both $S$ and $T$.
    Since each edge in $F$ is incident to precisely $k$ intervals $I_x$, $x \in \cX$, we have
    \begin{equation*}
        \sum_{x \in \cX} e_F(A,I_x) = k e(F) = k M.
    \end{equation*}
    Consider the following set of pairs
    \begin{equation*}
        \cP \defined \set[\bigg]{(x, \set{e,e'}) \in \cX \times \binom{E(F)}{2} \st \text{$e$ and $e'$ are incident to $I_x$}}.
    \end{equation*}
    We double count the size of $\cP$.
    First, we provide a lower bound on $\cP$.
    Note that
    \begin{equation*}
        \card{\cP} = \sum_{x \in \cX} \binom{e_F(A,I_x)}{2} \geq N' \binom{\frac{1}{N'}\sum_{x \in \cX} e_F(A,I_x)}{2} = N' \binom{k M / N'}{2},
    \end{equation*}
    by Jensen's inequality and the convexity of $t \mapsto \binom{t}{2}$.

    Now we give an upper bound on $\cP$ in terms of the number of edges in $\cH_k(A,U)$.
    Consider a pair $(x,\set{e,e'}) \in \cP$ where the edges $e$ and $e'$ are of the form $e = \set{a, a + u}$ and $e' = \set{a', a' + u'}$, for uniquely determined $a, a' \in A$ and $u, u' \in U$, and we assume $a \leq a'$.
    We split the pairs in $\cP$ into three classes $\cP_1$, $\cP_2$, and $\cP_3$.

    \begin{description}
        \item[$\cP_1$, the pairs with $a < a'$ and $u \neq u'$]
              these corresponds precisely to edges in $\cH_k(A,U)$, where the edge between $a$ and $a'$ and of type $(u, u')$ is counted with multiplicity $k - r$.
              Therefore, we have
              \begin{align*}
                  \card{\cP_1} & = \sum_{\substack{u, u' \in U, \, u \neq u' \\ a, a' \in A, \, a < a'}} \card{\set{x \in \cX \st a + u, a' + u' \in I_x}} \\
                               & = \sum_{\substack{u, u' \in U, \, u \neq u' \\ a, a' \in A, \, a < a'}} \max\set{0,k-\abs{(a + u) - (a' + u')}}
                  \leq k\, e(\cH_k(A,U)).
              \end{align*}

        \item[$\cP_2$, the pairs with $a < a'$ and $u = u'$]
              We have
              \begin{align*}
                  \card{\cP_2} & = \sum_{\substack{u \in U                                                                   \\ a, a' \in A, \, a < a'}} \card{\set{x \in \cX \st a + u, a' + u \in I_x}} \\
                               & = \sum_{r = 1}^\infty \sum_{u \in U} \sum_{a \in A} \ind\set{a + r \in A} \max\set{0,k - r} \\
                               & \leq \sum_{r = 1}^k \sum_{u \in U} \sum_{a \in A} k = k^2 M.
              \end{align*}

        \item[$\cP_3$, the pairs with $a = a'$]
              When $a = a'$ then we must have $u \neq u'$ and furthermore, we may assume that $u < u'$.
              We have
              \begin{align*}
                  \card{\cP_3} & = \sum_{\substack{u, u' \in U, \, u < u'                                                    \\ a \in A}} \card{\set{x \in \cX \st a + u, a + u' \in I_x}} \\
                               & = \sum_{r = 1}^\infty \sum_{a \in A} \sum_{u \in U} \ind\set{u + r \in U} \max\set{0,k - r} \\
                               & \leq \sum_{r = 1}^k \sum_{a \in A} \sum_{u \in U} k = k^2 M.
              \end{align*}
    \end{description}

    Therefore, we have
    \begin{equation*}
        \card{\cP} = \card{\cP_1} + \card{\cP_2} + \card{\cP_3} \leq k\, e(\cH_k(A,U)) + 2 k^2 M.
    \end{equation*}

    Combining the lower and upper bounds on $\card{\cP}$, we obtain
    \begin{align*}
        k\, e(\cH_k(A,U)) & \geq N' \binom{kM/N'}{2} - 2 k^2 M                               \\
                          & = N' \cdot \frac{kM}{2N'} \p[\Big]{\frac{kM}{N'} - 1} - 2 k^2 M.
    \end{align*}
    Since $N \geq k$, we have $N' \leq 3N$.
    Thus
    \begin{align*}
        e(\cH_k(A,U)) & \geq \frac{M}{2} \p[\Big]{\frac{kM}{N'} - 1 - 4k} \\
                      & \geq \frac{M}{2} \p[\Big]{\frac{kM}{3N} - 1 - 4k}
        \geq \frac{kM^2}{8N} = \frac{k \card{A}^2 \card{U}^2}{8N},
    \end{align*}
    where we have used that $kM/3N - 4k - 1 \geq kM/4N$, which follows from the hypothesis that $M \geq 60 N$.
\end{proof}
%--------------------------------------%




%--------------------------------------%
%--------------------------------------%
\begin{proposition}
    \label{prop:graph-containers}
    Let \(G\) be a graph on \(n\) vertices, \(q,r\in\mathbb Z_{\geq0}\), \(0<\beta<1\), and \(R\in\mathbb Z_{\geq0}\) with \(R\geq e^{-\beta q}n\).
    Suppose that every \(S\subseteq V(G)\) with \(\card{S}\geq R\) satisfies
    \[
        e(G[S])
        \geq
        \beta\binom{\card{S}}{2}.
    \]
    Then the number of independent sets of \(G\) of cardinality
    \(q+r\) is at most \(\binom nq\binom Rr\).
\end{proposition}
%--------------------------------------%

\begin{proposition}
    \label{prop:containers-sidon}
    Let \(k,u,N\in\mathbb Z_{\geq1}\) satisfy \(u\geq32\) and \(N\geq ku^2\).
    Let \(q,r\in\mathbb Z_{\geq0}\) satisfy
    \begin{equation}
        \label{eq:cont-condition}
        \frac{ku^2q}{16N}
        \geq
        \log\left(\frac{u}{64}\right).
    \end{equation}
    Define \(R\defined\ceil{\frac{64N}{u}}\).
    Then
    \[
        S_k(N,u+q+r)
        \leq
        S_k(N,u)
        \binom Nq
        \binom Rr.
    \]
\end{proposition}
%--------------------------------------%

\begin{proof}
    Fix a \(k\)-strong Sidon set \(U\subseteq[N]\) with \(\card{U}=u\),
    and set \(A\defined[N]\setminus U\). We bound the number of \(k\)-strong Sidon sets \(T\subseteq[N]\)
    with \(U\subseteq T\) and \(\card{T}=u+q+r\).
    Writing \(T=U\cup I\) with \(I\subseteq A\) and \(\card{I}=q+r\),
    Proposition~\ref{prop:sidon-independent} implies that \(I\) is
    an independent set of \(\cH_k(A,U)\). By Lemma~\ref{lem:k-separated-core}, there exists \(U'\subseteq U\)
    with \(\card{U'}\geq u-1\) such that \(U'\) is \(k\)-separated.
    Since the \(k\)-strong Sidon property is hereditary, \(U'\) is
    also \(k\)-strong Sidon. Since \(U'\subseteq U\), the multigraph \(\cH_k(A,U')\)
    is a spanning submultigraph of \(\cH_k(A,U)\), so every
    independent set of \(\cH_k(A,U)\) is also independent in
    \(\cH_k(A,U')\). Let \(G\) be the underlying simple graph of
    \(\cH_k(A,U')\). By Corollary~\ref{cor:multiplicity}, every edge
    of \(\cH_k(A,U')\) has multiplicity at most \(2\).
    Define
    \[
        \beta\defined\frac{ku^2}{16N}.
    \]
    Since \(N\geq ku^2\), we have \(0<\beta\leq\frac1{16}<1\).
    By \eqref{eq:cont-condition}, \(\beta q\geq\log(u/64)\),
    and therefore \(e^{-\beta q}\leq 64/u\).
    Since \(G\) has \(\card{A}=N-u\) vertices,
    \[
        e^{-\beta q}\card{A}
        \leq
        e^{-\beta q}N
        \leq
        \frac{64N}{u}
        \leq
        R.
    \]

    It remains to verify the local density condition.
    Let \(S\subseteq A\) with \(\card{S}\geq R\).
    Since \(\card{U'}\geq u-1\),
    we have
    \begin{align*}
        \card{S}\card{U'}
         & \geq
        R(u-1)                    \\
         & \geq
        \frac{64N}{u}(u-1)        \\
         & =
        64N\left(1-\frac1u\right) \\
         & \geq
        60N,
    \end{align*}
    where the last inequality follows from \(u\geq32\).

    Hence Lemma~\ref{lem:supersaturation} applies to \(S\) and \(U'\),
    and gives
    \[
        e\bigl(\cH_k(S,U')\bigr)
        \geq
        \frac{k\card{S}^2\card{U'}^2}{8N}
        \geq
        \frac{k\card{S}^2(u-1)^2}{8N}.
    \]

    Since every edge of \(\cH_k(S,U')\) has multiplicity at most \(2\),
    \(e(G[S])\geq\frac12 e\bigl(\cH_k(S,U')\bigr)\),
    and therefore \(e(G[S])\geq\frac{k\card{S}^2(u-1)^2}{16N}\).
    Since \(u\geq32\), we have \(u^2\leq2(u-1)^2\).
    Thus
    \begin{align*}
        e(G[S])
         & \geq
        \frac{k\card{S}^2(u-1)^2}{16N} \\
         & \geq
        \frac{ku^2\card{S}^2}{32N}     \\
         & =
        \frac{\beta\card{S}^2}{2}      \\
         & \geq
        \beta\binom{\card{S}}{2}.
    \end{align*}

    Hence all the hypotheses of
    Proposition~\ref{prop:graph-containers} are satisfied. Therefore
    the number of independent sets of \(G\) of cardinality \(q+r\) is
    at most
    \[
        \binom{N-u}{q}\binom Rr
        \leq
        \binom Nq\binom Rr.
    \]

    Since every \(k\)-strong Sidon extension \(T\) of \(U\) gives such
    an independent set, the result follows.
\end{proof}
%-----------------------------------






\end{document}